\documentclass[letterpaper,10pt]{article}

\usepackage{latexsym}
\usepackage[empty]{fullpage}
\usepackage{titlesec}
\usepackage{marvosym}
\usepackage[usenames,dvipsnames]{color}
\usepackage{verbatim}
\usepackage{enumitem}
\usepackage[hidelinks]{hyperref}
\usepackage{fancyhdr}
\usepackage{tabularx}
\usepackage{multicol}
\usepackage[utf8]{inputenc}
\usepackage{CJKutf8}
\input{glyphtounicode}

\pagestyle{fancy}
\fancyhf{}
\fancyfoot{}
\renewcommand{\headrulewidth}{0pt}
\renewcommand{\footrulewidth}{0pt}


\addtolength{\oddsidemargin}{-0.5in}
\addtolength{\evensidemargin}{-0.5in}
\addtolength{\textwidth}{1in}
\addtolength{\topmargin}{-.5in}
\addtolength{\textheight}{1.0in}

\urlstyle{same}

\raggedbottom
\raggedright
\setlength{\tabcolsep}{0in}

\titleformat{\section}{
  \vspace{-4pt}\centering
}{}{0em}{}[\color{black}\titlerule\vspace{-5pt}]


\pdfgentounicode=1

\newcommand{\resumeItem}[1]{
  \item\small{
    {#1 \vspace{-2pt}}
  }
}

\newcommand{\resumeSubheading}[4]{
  \vspace{-2pt}\item
    \begin{tabular*}{0.97\textwidth}[t]{l@{\extracolsep{\fill}}r}
      \textbf{#1} & #2 \\
      \textit{\small#3} & \textit{\small #4} \\
    \end{tabular*}\vspace{5pt}
}

\newcommand{\resumeSubSubheading}[2]{
    \item
    \begin{tabular*}{0.97\textwidth}{l@{\extracolsep{\fill}}r}
      \textit{\small#1} & \textit{\small #2} \\
    \end{tabular*}\vspace{-7pt}
}

\newcommand{\resumeProjectHeading}[2]{
    \item
    \begin{tabular*}{0.97\textwidth}{l@{\extracolsep{\fill}}r}
      \small#1 & #2 \\
    \end{tabular*}\vspace{-7pt}
}

\newcommand{\resumeSubItem}[1]{\resumeItem{#1}\vspace{-4pt}}

\renewcommand\labelitemii{$\vcenter{\hbox{\tiny$\bullet$}}$}

\newcommand{\resumeSubHeadingListStart}{\begin{itemize}[leftmargin=0.15in, label={}]}
\newcommand{\resumeSubHeadingListEnd}{\end{itemize}}
\newcommand{\resumeItemListStart}{\begin{itemize}}
\newcommand{\resumeItemListEnd}{\end{itemize}\vspace{-5pt}}
\newcommand{\resumeDesc}[1]{%
  \vspace{-4pt}%
  \begin{flushleft}\small #1\end{flushleft}%
  \vspace{-6pt}%
}
\newcommand{\Name}[1]{\fontsize{28}{30}\selectfont #1}

\begin{document}
\begin{CJK}{UTF8}{gbsn}



\begin{center}
    {\Name James Fang}  \\ \vspace{-2pt}
    \begin{multicols}{2}
    \begin{flushleft}
    \href{{https://github.com/YaYa-xhy}}{Github: https://github.com/yayaxhy}\\
    \end{flushleft}
    
    \begin{flushright}
    \href{mailto:jamesfang1116@gmail.com}{Email: jamesfang1116@gmail.com}
    \item{电话: 437-662-9197}
    \end{flushright}
    \end{multicols}
\end{center}


%-----------EDUCATION-----------
\vspace{-20pt}
\section{教育背景}
  \resumeSubHeadingListStart
      \resumeSubheading
      {多伦多都市大学}{2023年9月 -- 至今}
      {计算机科学本科三年级}{多伦多, 加拿大}

  \resumeSubHeadingListEnd
      



%-----------EXPERIENCE-----------
\section{经历}
 \resumeSubHeadingListStart
 \resumeSubheading 
      {Discord 客服}{2023年3月 -- 至今}
      {} {远程}\\ 
      \\resumeItemListStart
        \\resumeItem{在 1.7 万+ 成员的 Discord 服务器远程工作，负责问题排查与即时答疑。}
        \\resumeItem{维护用户关系，沉淀常见问题处理流程，提升处理效率。}
        \\resumeItem{每两个月组织并管理一次游戏活动，促进社区活跃度。}
      \\resumeItemListEnd
      
  \resumeSubHeadingListEnd
  
 \resumeSubHeadingListStart
 \resumeSubheading 
      {Discord 技术负责人（初创游戏社群）}{2024年12月}
      {} {远程}\\ 
      \\resumeItemListStart
        \\resumeItem{负责 Discord 机器人与关联网页的开发、维护与功能创新。}
        \\resumeItem{根据运营需求迭代点单、打赏、榜单等核心流程。}
        \\resumeItem{跟进线上问题与数据异常，保障功能稳定运行。}
      \\resumeItemListEnd
  \resumeSubHeadingListEnd



%-----------PROGRAMMING SKILLS-----------
\section{技术技能、语言与兴趣}
 \begin{itemize}[leftmargin=0.15in, label={}]
    \small{\item{
    \textbf{操作系统}{: Windows, Linux} \\
     \textbf{编程语言}{: Python, Java, SQL} \\
     \textbf{办公}{: Office} \\
     \textbf{语言}{: 英语 (IELTS 7.5), 中文 (母语)} \\
     \textbf{兴趣}{: K-pop}
     
    }}
 \end{itemize}

%-----------CERTIFICATIONS-----------
\section{项目}
  \resumeSubHeadingListStart
 \resumeSubheading 
      {Discord 陪玩机器人与管理后台/礼物墙网站}{}
      {} {}\\
      \resumeItemListStart
        \resumeItem{从 0 到 1 搭建点单、打赏、排行榜、礼物墙、积分、结算与通知功能。}
        \resumeItem{设计可审计的交易流水，支持回滚与幂等处理，保障数据一致性。}
        \resumeItem{实现按分类集齐礼物墙奖励并自动通知用户。}
        \resumeItem{运行于 1 个服务器，用户 1000+，日均处理 50+ 订单。}
        \resumeItem{技术栈：TypeScript, Node.js, Prisma, PostgreSQL, Discord.js, Next.js。}
      \resumeItemListEnd
  \resumeSubHeadingListEnd

  \resumeSubHeadingListStart
 \resumeSubheading 
      {Discord 音乐机器人}{}
      {} {}\\ 
      可接收 YouTube 链接并播放音频。\\
      
  \resumeSubHeadingListEnd

\end{CJK}
\end{document}
